\documentclass[11pt]{article}
\usepackage[margin=1in]{geometry}
\usepackage{amsmath,amssymb,amsthm}

\newtheorem{theorem}{Theorem}
\newtheorem{remark}{Remark}

\title{Phi-Descent Closure with Explicit Topological Input}
\author{}
\date{}

\begin{document}
\maketitle

\begin{theorem}[Conditional $\Phi$-closure]
Let $M$ be a closed simply connected smooth $3$-manifold and let
\[
\Phi(g):=-\lambda_1(-4\Delta_g+s_g).
\]
Assume:
\begin{enumerate}
\item $\Phi$ is non-increasing along the chosen normalized Ricci-flow-with-surgery trajectory;
\item $\Phi(g_\ast)=0$ at a limit metric $g_\ast$;
\item vanishing of the monotonicity defect yields the gradient shrinking soliton equation
\[
\operatorname{Ric}_{g_\ast}+\nabla^2 f=\mu g_\ast
\qquad\text{with}\qquad \mu>0.
\]
\end{enumerate}
Then $M$ is diffeomorphic to $S^3$.
\end{theorem}

\begin{proof}
From the shrinking soliton equation with $\mu>0$, tracing gives positive average scalar curvature.
On a compact shrinking soliton this implies positive scalar curvature somewhere, hence everywhere after
standard maximum-principle propagation along the normalized flow.
In dimension $3$, an Einstein metric has constant sectional curvature.
Thus the limiting metric is positively curved.
Bonnet--Myers implies finite fundamental group; since $M$ is simply connected, the group is trivial.
A closed simply connected $3$-manifold with constant positive sectional curvature is homeomorphic to
$S^3$.
By Moise's theorem, every topological $3$-manifold admits a unique smooth structure up to diffeomorphism,
so $M$ is diffeomorphic to $S^3$.
\end{proof}

\begin{remark}
This closes the descent only after the soliton/positive-curvature step is established.
The remaining analytic burden is the derivation of the shrinking-soliton conclusion directly from the
$\Phi$-framework.
\end{remark}

\end{document}
